\chapter{LOGO}\label{cha:logo}

LOGO is the language that taught a generation of children that a
computer does what it is told, and it did it with a turtle: a small
arrow on the screen that walks forward, turns, and leaves a line behind
it. The K4510's is new, written for this machine in September~2026, and
native --- 45GS10 code, no co-processor.

\begin{type}
LOGO
FD 100 RT 90
REPEAT 4 [FD 100 RT 90]
\end{type}

\screen{shots/logo.png}{LOGO: a square repeated round a circle, drawn
by the blitter under the console's text.}

The picture is VICKY's 640$\times$480 bitmap \emph{under} the console,
which is the classic LOGO screen: the turtle draws below and the prompt
lives above, so you can see what you asked for and what you are asking
next at once. The turtle itself is a small green turtle, a hardware
sprite in sixteen headings --- one every 22.5 degrees, drawn once and
turned by \pth{tools/mkturtle.py} when the machine is built, so turning
it is a matter of pointing the sprite at the right picture. LOGO puts the
console into 640$\times$480 for the session, as EhBASIC's
\verb|GRAPHICS 2| does, and back into the mode it found at \verb|BYE|.
Every line is the
blitter's, and every number is a real one --- IEEE floating point, done
by the MATH unit, so a heading is an angle and \verb|SQRT 2| is what it
should be.

\section{The words}
\begin{center}
\small
\begin{tabular}{@{}>{\raggedright\arraybackslash}p{0.42\linewidth} >{\raggedright\arraybackslash}p{0.52\linewidth}@{}}
\toprule
\texttt{FD n}, \texttt{BK n} & forward, back \\
\texttt{RT n}, \texttt{LT n} & right, left, in degrees \\
\texttt{PU}, \texttt{PD} & pen up, pen down \\
\texttt{HT}, \texttt{ST} & hide, show the turtle \\
\texttt{HOME}, \texttt{CS} & to the centre; clear the screen as well \\
\texttt{SETXY x y} & put it somewhere \\
\texttt{SETX}, \texttt{SETY}, \texttt{SETH} & one coordinate, or the heading \\
\texttt{SETPC n} & the pen's colour \\
\texttt{FILL} & paint the area under the turtle \\
\texttt{XCOR}, \texttt{YCOR}, \texttt{HEADING} & where it is \\
\texttt{REPEAT n [\ldots]} & do it n times; \texttt{REPCOUNT} says which time this is \\
\texttt{IF c [\ldots]} & decide \\
\texttt{IFELSE c [\ldots] [\ldots]} & decide between two \\
\texttt{MAKE "x expr}, \texttt{:x} & a variable, and its value \\
\texttt{PRINT}, \texttt{SHOW} & print a value \\
\texttt{RANDOM}, \texttt{SQRT}, \texttt{SIN}, \texttt{COS}, \texttt{INT}, \texttt{ABS} & numbers \\
\texttt{TO name :a \ldots{} END} & a procedure; \texttt{STOP} and \texttt{OUTPUT} leave it \\
\texttt{LOAD "name} & run a \texttt{.LGO} file \\
\texttt{EDIT "name} & edit one in VI, then run it \\
\texttt{HELP}, \texttt{BYE} & the words; back to K/OS \\
\bottomrule
\end{tabular}
\end{center}

\section{Procedures}
A procedure is a new word, and it may call itself:

\begin{type}
TO TREE :N
  IF :N < 5 [STOP]
  FD :N LT 30 TREE :N * 0.7
  RT 60 TREE :N * 0.7
  LT 30 BK :N
END
TREE 80
\end{type}

\noindent A procedure's inputs are its own; \verb|MAKE| writes the
nearest variable of that name, or makes a global one. Recursion is real
and goes deep --- deep enough for any tree --- and a procedure that
recurses for ever says so rather than taking the machine with it.

\section{Files}
\verb|LOAD "TREE| runs \texttt{TREE.LGO} from the directory you are in,
or from the examples in \pth{/LANG/LOGO/EX}: \texttt{SQUARE},
\texttt{SPIRAL}, \texttt{TREE} and \texttt{SNOW}. \verb|EDIT "MINE|
opens \texttt{MINE.LGO} in VI (Chapter~\ref{cha:editors}) and runs it
when you leave the editor, which is the way to write anything longer
than a line. A \texttt{.LGO} file is plain text, one line as you would
type it after another.
